\documentclass{ximera}

\input{../../preamble.tex}

\title{Exponential Analysis}

\begin{document}

\begin{abstract}
%Stuff can go here later if we want!
\end{abstract}
\maketitle




Precalculus is largely about the structure of the Elementary Functions and how we reason about them.  This includes power functions; radical and root functions; linear, quadratic, and other polynomial functions; rational functions; exponential functions; logarithmic functions; and a library of trigonometric functions (which are introduced in this course).


Our focus is analyzing these functions.





\begin{itemize}
     \item \textbf{\textcolor{red!70!black}{Domain}} 
     \item \textbf{\textcolor{red!70!black}{Zeros}} 
     \item \textbf{\textcolor{red!70!black}{Continuity}} \\
     \text{    }$\circ$ \textbf{\textcolor{purple!85!blue}{discontinuities}}\\
     \text{    }$\circ$ \textbf{\textcolor{purple!85!blue}{singularities}}\\
     \item \textbf{\textcolor{red!70!black}{End-Behavior}} 
     \item \textbf{\textcolor{red!70!black}{Behavior}} \\
     \text{    }$\circ$ \textbf{\textcolor{purple!85!blue}{intervals where increasing}}\\
     \text{    }$\circ$ \textbf{\textcolor{purple!85!blue}{intervals where decreasing}}\\
     \item \textbf{\textcolor{red!70!black}{Global Maximum and Minimum}} 
     \item \textbf{\textcolor{red!70!black}{Local Maximums and Minimums}} 
     \item \textbf{\textcolor{red!70!black}{Range}} 
     \item \textbf{\textcolor{blue!55!black}{...and we would like a nice graph}} 
\end{itemize}






In this section, we'll analyze exponential functions.




$\blacktriangleright$ \textbf{\textcolor{blue!55!black}{Core Exponential Functions}} 

The template for the core exponential function looks like \textbf{\textcolor{blue!55!black}{$exp(x) = r^x$}}

It is a decreasing function, if $r < 1$.\\
It is an increasing function, if $r > 1$.


Composiing the core exponential functions with linear functions give us our ELementary Functions version of exponential functions.


$\blacktriangleright$ \textbf{\textcolor{blue!55!black}{Exponential Functions}} 


The template for the basic exponential function looks like \textbf{\textcolor{blue!55!black}{$exp(x) = A \cdot r^x$}} (\textbf{\textcolor{blue!55!black}{$exp(x) = A \cdot r^{B \, x + C}$}}). 

\begin{itemize}
\item $A$ is the \textbf{\textcolor{purple!85!blue}{leading coefficient}}.  It primarily determines the sign of the function values.
\item $r$ is the \textbf{\textcolor{purple!85!blue}{base}}. It determines how the function values grow, depending on whether $0 < r < 1$ or $1 < r$.
\end{itemize}





$\blacktriangleright$ \textbf{\textcolor{blue!55!black}{Shifted Exponential Functions}}   

The template for a shifted exponential function looks like \textbf{\textcolor{blue!55!black}{$shexp(x) = A \cdot r^x + D$}}  (\textbf{\textcolor{blue!55!black}{$A \cdot r^{B \, x + C} + D$}}). 

\begin{itemize}
\item $C$ describes a \textbf{\textcolor{purple!85!blue}{shift}} in the domain from the basic exponential function.  
\item $D$ describes a \textbf{\textcolor{purple!85!blue}{shift}} in the function value from the basic exponential function. 
\end{itemize}

The first term in the formula for a shifted exponential function is, by itself, an exponential function:

\[
a \cdot r^{b \, x + c} = a \cdot r^{b \, x} \cdot  r^c = (a \cdot r^c) \cdot r^{b \, x} = (a \cdot r^c) \cdot (r^b)^x 
\]

It \textbf{\textcolor{purple!85!blue}{CAN}} be written in the form $A \cdot R^x$, where $A = a \cdot r^c$ and $R = r^b$.


Algebraically, the constant term, the ``+ D'' cannot be absorbed into the exponential form.  

Adding a constant onto an exponential function breaks the constant percentage growth.

That is why we have two exponential categories.





\subsection*{Learning Outcomes}


\begin{sectionOutcomes}
In this section, students will 

\begin{itemize}
\item describe the full picture of all (shifted) exponential functions and their graphs.
\end{itemize}
\end{sectionOutcomes}












\begin{center}
\textbf{\textcolor{green!50!black}{ooooo-=-=-=-ooOoo-=-=-=-ooooo}} \\

more examples can be found by following this link\\ \link[More Examples of Exponential Functions]{https://ximera.osu.edu/csccmathematics/precalculus2/precalculus2/expFunctions/examples/exampleList}

\end{center}









\end{document}
